\chapter{Gated losses, block masking, the decoder bottleneck and FT09}\label{ch:decoderwork}

\begin{onething}
Changing what the training loss sees (gating the worst pixels, hiding tokens in blocks) cost a little
detection, because most of what a model gets wrong is ordinary hard background; a diagnosis without training
then showed that FT04's error is mostly detail its decoder cannot render, because each $4\times4$ token is
squeezed into 32 numbers with no full-resolution path back, and the progressive decoder built to fix that
(FT09) raised contrast and gave a gain that cleared the bootstrap by a hair and then did not survive five more
epochs.
\end{onething}

This chapter covers three days of experiments, 2026-10-21 to 2026-10-23, after the uncertainty head of
\cref{ch:uncertainty} failed. It is three stories with one thread. The first two runs (FT07, FT08) changed the
\emph{training signal}. The diagnosis that followed (``is the decoder the bottleneck?'') asked where the error
actually comes from. The third run (FT09) changed the \emph{shape of the network}. Each story carries a piece
of mathematics: quantiles and what a gate removes from a gradient; expected numbers of visible neighbours under
random masking; projection onto a span and the Pythagorean split of an error; a counting argument about
information through a bottleneck; and a by-hand parameter count of a new decoder.

Provenance tags abbreviate paths as in \cref{ch:uncertainty}: \texttt{FT07-gated/}, \texttt{FT09-decoder/} and
the like are folders under \file{research/experiments/indradhanu-finetune/}; \texttt{decoder-capacity/} and
\texttt{bench-resolution/} are under \file{research/experiments/}.

\begin{analogy}
A police sketch artist draws faces from a witness's description, and the sketches come out poor. Is the
witness's memory the problem, or the artist's pencil? You could coach the witness to leave out the details
they remember worst (that is FT07), or to describe the face in bigger pieces (FT08), and the sketches barely
change. The decisive test is to put the real face in front of the artist and ask for a copy. If the copy is
barely better than the sketch from memory, the memory was never the limit: the pencil is. That test is
\cref{sec:decoderwork:diagnosis}, and FT09 is a finer pencil.
\end{analogy}

%======================================================================
\section{The question behind FT07 and FT08}\label{sec:decoderwork:question}

A masked autoencoder fine-tuned on shards that contain real targets can learn to reconstruct them, and a
target it reconstructs well leaves no residual to score (\cref{ch:recon}). GT-HAD \citep{lian2025gated}
\emph{gates} its training: it keeps likely anomalies out of what the network learns from. The anomaly
course tried a gate per scene at inference time and found that a short run never learns the targets anyway,
so the gate changed nothing (\cref{ch:scoringgains}). Five epochs over the whole training mix is different: it
sees every target-like pixel many times. So two runs, both from FT04's \code{best.pt}, 5 epochs, learning rate
$5\times10^{-5}$ on a cosine schedule, FT04's configuration otherwise
\tagM{FT07-gated/README.md}:

\begin{center}
\begin{tabular}{lll}
\toprule
run & change & configuration fields\\
\midrule
FT07 & gated loss & \texttt{gate\_pct: 1.0}, \texttt{gate\_halo: 2}\\
FT08 & gated loss + $2\times2$-token block masking & as FT07, plus \texttt{mask\_block: 2}\\
\bottomrule
\end{tabular}
\end{center}

The expectation, written before launch: FT07 a small gain at the top 0.1\% and slightly worse reconstruction;
FT08 the larger detection change, in either direction.

%======================================================================
\section{Gating the worst 1\%}\label{sec:decoderwork:gate}

\subsection{Quantiles}

\begin{definition}[Quantile]
The $q$-\emph{quantile}\index{quantile} of a list of numbers, for $0\le q\le 1$, is the value below which a
fraction $q$ of the list lies. With $n$ sorted values $x_{(0)}\le x_{(1)}\le\dots\le x_{(n-1)}$ (counting from
0), \code{torch.quantile} uses the position $q(n-1)$ and interpolates linearly between the two neighbouring
sorted values: at position $k+f$ ($k$ an integer, $0\le f<1$) the quantile is $x_{(k)}+f\,(x_{(k+1)}-x_{(k)})$.
\end{definition}

The function \code{gate_pixels} in \file{app/foundation_models/trainers/finetune_trainer.py} does four
things, for each source group (each sensor's tiles) in a mixed batch (\cref{ch:finetune}):
\begin{lstlisting}
err  = per-pixel L1, mean over valid bands                        # (B, 1, H, W)
cut  = torch.quantile(err[loss pixels][:1_000_000], 1 - pct/100)  # the group's top pct % threshold
gate = (err >= cut) & loss_pixel
gate = max_pool2d(gate, 2*halo + 1, stride=1, padding=halo) > 0   # dilate to a (2 halo + 1)-wide square
\end{lstlisting}
and the reconstruction loss mask becomes \code{loss_mask * ~gate}. Three details matter.
\begin{itemize}
\item \code{x_hat} is \emph{detached} before the gate is computed: the gate is a selection made without
  gradient, so the model cannot learn to change which pixels get gated.
\item The quantile is taken over \emph{loss pixels only} (hidden, valid pixels), since only they carry loss.
\item At most the first 1,000,000 loss pixels enter the quantile, because \code{torch.quantile} rejects very
  large inputs. Eight 128\,px tiles are $8\times128\times128=131{,}072$ pixels, so at FT07's batch size of 8
  the cap is never reached.
\end{itemize}

\begin{worked}[label={we:decoderwork:quantile}]{The top 1\% of 100 loss pixels, with its halo}
Do this on paper. A group has 100 loss pixels laid out as a $10\times10$ patch, all hidden, with distinct
errors.
\begin{align*}
\text{quantile level} &= 1-1/100 = 0.99\\
\text{position} &= 0.99\times(100-1) = 98.01 \quad(k=98,\ f=0.01)\\
\text{cut} &= x_{(98)} + 0.01\,(x_{(99)}-x_{(98)})\\
&\phantom{=}\ \text{(1\% of the way from the 99th to the 100th smallest)}\\
\text{err}\ge\text{cut} &: \text{only } x_{(99)}, \text{the single largest error}\\
\text{gated before the halo} &= 1 \text{ pixel} = 1\% \text{ of } 100
\end{align*}
Now the halo, \code{gate_halo = 2}: \code{max_pool2d} with kernel $2\times2+1=5$ turns one pixel into a
$5\times5$ square.
\begin{align*}
\text{gated pixel in the middle} &: 5\times5 = 25 \text{ pixels};\quad 100-25 = 75 \text{ loss pixels left}\\
\text{gated pixel in a corner} &: 3\times3 = 9 \text{ pixels};\quad 100-9 = 91 \text{ loss pixels left}
\end{align*}
A 1\% gate removed a quarter of this group's loss pixels.
\end{worked}

\subsection{What fraction of the gradient a gate removes}\label{sec:decoderwork:gradfrac}

``Removing 1\% of the pixels'' sounds small. Two different things could be meant by how much is removed: the
share of the \emph{loss value}, or the share of the \emph{gradient}. For L1 they differ sharply.

The L1 loss over $N$ loss pixels and $C$ bands is $L=\frac{1}{NC}\sum_{p,c}\abs{\hat x_{pc}-x_{pc}}$.
Differentiate with respect to one reconstructed value $\hat x_{pc}$. Only one term contains it, and the
derivative of $\abs{u}$ is $\sign(u)$ (for $u\ne0$):
\begin{equation}\label{eq:decoderwork:l1grad}
\frac{\partial L}{\partial\hat x_{pc}} = \frac{1}{NC}\,\sign(\hat x_{pc}-x_{pc}).
\end{equation}
Every included pixel-band pushes on the network with the \emph{same strength}, $1/(NC)$, however large its
error. So the share of gradient terms the gate removes is the share of loss pixels it removes, while the share
of loss \emph{value} it removes is larger, because the gated pixels are the large ones. (Under an L2 loss the
gradient would be proportional to the error, and the worst pixels would dominate the gradient as well.)

\begin{worked}[label={we:decoderwork:share}]{Loss share against gradient share}
A group of 100 loss pixels: 99 have error $0.05$ and one has error $0.50$ (illustrative numbers \tagA). Gate
the worst one, without a halo.
\begin{align*}
\text{total loss value} &= 99\times0.05 + 0.50 = 4.95 + 0.50 = 5.45\\
\text{share of loss value removed} &= 0.50/5.45 = 0.092 \quad(9.2\%)\\
\text{share of L1 gradient terms removed} &= 1/100 = 0.01 \quad(1\%)\\
\text{under L2: } \textstyle\sum 2\abs{r} &= 99\times0.10 + 1.00 = 9.90+1.00 = 10.90\\
\text{share of L2 gradient size removed} &= 1.00/10.90 = 0.092 \quad(9.2\%)
\end{align*}
With the halo of 25 pixels, the gradient share is 25\%, and almost all of those 24 extra pixels are ordinary.
\end{worked}

How large the gated share was in FT07 is not recorded: the run logged no gated share. You can bound it.

\begin{worked}[label={we:decoderwork:tile}]{An upper bound on a real batch}
Use FT07's batch of 8 tiles of $128\times128$ and assume (\tagA) about 65\% of pixels are hidden and valid, and
the gated pixels are isolated, far from edges and from each other.
\begin{align*}
\text{pixels} &= 8\times128\times128 = 131{,}072\\
\text{loss pixels} &\approx 0.65\times131{,}072 = 85{,}197\\
\text{gated before halo} &\approx 0.01\times85{,}197 = 852\\
\text{loss pixels per square} &\approx 0.65\times25 = 16.25\quad\text{(the square also covers visible pixels)}\\
\text{removed} &\approx 852\times16.25 = 13{,}845\\
\text{share} &\approx 13{,}845/85{,}197 = 0.16
\end{align*}
At most about 16\%. Real high-error pixels come in clumps whose squares overlap, so the true share is lower.
\end{worked}

\section{Block masking}\label{sec:decoderwork:blocks}

\subsection{Why blocks}

\Cref{ch:uncertainty} found that the head could see a target ``through the neighbours'': a target's halo
reaches into neighbouring tokens, and a model filling a hidden token can copy the odd spectrum from next door.
Hide bigger units and less of the target is left in view. FT08 hides tokens in $2\times2$ blocks, 8\,px on a
side.

\begin{worked}[label={we:decoderwork:footprint}]{A 5\,px footprint on a 4\,px and an 8\,px grid}
A one-pixel target whose light reaches 2\,px either side has a footprint $2+1+2=5$\,px wide. Assume (\tagA) its
starting offset on the grid is uniform.
\begin{align*}
\text{4\,px token} &: 5>4,\ \text{never fits}: 0 \text{ of } 4 \text{ offsets}\\
\text{8\,px block} &: \text{fits if it starts at offset } 0,1,2,3:\ 4 \text{ of } 8 = 1/2 \text{ per axis}\\
\text{both axes} &: 1/2\times1/2 = 1/4
\end{align*}
So a block contains the whole footprint about a quarter of the time; otherwise it straddles a border and
leaks unless the neighbouring block is hidden too.
\end{worked}

\subsection{How the block mask is drawn}

\code{_build_prediction_mask} in the fine-tune trainer, with \code{mask_block = b > 1}:
\begin{lstlisting}
tokens = valid-token mask                                  # (B, N), N = Ht * Wt
grid   = tokens.reshape(B, 1, Ht, Wt)
blocks = max_pool2d(grid, b, ceil_mode=True) > 0          # (B, 1, Ht/b, Wt/b): valid if any token is
_, block_pred = generate_prediction_mask(blocks, mask_ratio)   # hide mask_ratio of the valid blocks
pred   = interpolate(block_pred, scale_factor=b, mode="nearest")[..., :Ht, :Wt] * tokens
\end{lstlisting}

\begin{worked}[label={we:decoderwork:hidden}]{How much a 128\,px tile hides}
At \code{mask_ratio} $0.65$:
\begin{align*}
\text{tokens per side} &= 128/4 = 32;\quad \text{tokens} = 32\times32 = 1{,}024\\
\text{blocks per side} &= 32/2 = 16;\quad \text{blocks} = 16\times16 = 256\\
\text{blocks hidden} &= 256\times0.65 = 166.4 \to 166 \ (\text{truncated})\\
\text{tokens hidden} &= 166\times4 = 664;\quad 664/1{,}024 = 0.648\\
\text{single tokens} &: 1{,}024\times0.65 = 665.6 \to 665
\end{align*}
FT08 hides almost exactly as much as FT04; it hides it in bigger pieces.
\end{worked}

\subsection{Expected visible neighbours}\label{sec:decoderwork:neighbours}

How much context does a hidden token keep? Count its \emph{visible neighbours}: of the 8 tokens around it
(sides and corners), how many are visible. Ignore tile borders (\tagA).

\paragraph{Random single tokens.} Given that our token is hidden, the remaining $1{,}023$ tokens contain
$665-1=664$ hidden and $1{,}024-665=359$ visible ones. Each neighbour is visible with probability
$359/1{,}023$. By \emph{linearity of expectation}\index{linearity of expectation} (the expected count is the sum
of the probabilities, whether or not the events are independent):
\begin{equation}\label{eq:decoderwork:single}
\E[\text{visible neighbours}] = 8\times\frac{359}{1{,}023} = 8\times0.3509 = 2.807.
\end{equation}

\paragraph{$2\times2$ blocks.} Every token in a $2\times2$ block sits at a corner of its block. Take the
top-left token. Of its 8 neighbours, 3 are in its own block (hidden for certain), 2 in the block above, 2 in
the block to the left, 1 in the diagonal block.

\begin{figure}[htbp]
\centering
\begin{tikzpicture}[scale=0.7]
\fill[maskgrey!70] (2,2) rectangle (4,4);
\fill[loss!70] (2,3) rectangle (3,4);
\foreach \i/\j in {1/2,1/3,1/4,2/4,3/4} {\fill[spat!40] (\i,\j) rectangle ++(1,1);}
\foreach \x in {0,...,5} \foreach \y in {0,...,5} \draw[gray!50] (\x,\y) rectangle ++(1,1);
\foreach \x in {0,2,4} \foreach \y in {0,2,4} \draw[very thick] (\x,\y) rectangle ++(2,2);
\node[font=\small,align=left,anchor=west] at (6.3,4) {amber: the hidden token\\grey: its own block, hidden\\indigo: the 5 neighbours in\\\quad 3 other blocks};
\end{tikzpicture}
\caption{The top-left token of a hidden $2\times2$ block (thick outlines are blocks). Three of its 8 neighbours
are in its own block and are certainly hidden; the other 5 lie in three neighbouring blocks.}
\label{fig:decoderwork:neighbours}
\end{figure}

Given our block is hidden, the other 255 blocks contain $166-1=165$ hidden and $256-166=90$ visible ones. Each of
the 5 outside neighbours is visible with probability $90/255$:
\begin{align*}
\E[\text{visible neighbours}] &= 3\times0 + 5\times\frac{90}{255}\\
&= 5\times0.3529 = 1.765.
\end{align*}
Blocks cut the expected visible context from 2.807 to 1.765 neighbours, by a factor $1.765/2.807=0.63$. That is
the intended effect: less of a target's halo in view. The cost is also visible in the arithmetic: ordinary
texture is now rebuilt from neighbours further away.

\subsection{Inference must match}

The product reconstructs in two complementary checkerboard passes (\cref{ch:inference}). A model trained to
fill 8\,px holes but asked to fill 4\,px ones sees more context than it trained with. So FT08 had to be inferred
with checkerboard cells of \code{mask_block} tokens. On a $32\times32$-token tile with cell 2: $16\times16$ cells,
and each pass hides 128 cells $=512$ tokens $=50\%$.

The parent inferencer (\file{app/foundation_models/inferencers/segformer_mae_inferencer.py}) makes \emph{two}
checkerboard choices, both from \code{config.checkerboard_cell_size}: which tokens each pass hides (the keep masks,
built by \code{_checkerboard_keep_mask}) and which pass each pixel's prediction is taken from (\code{pass1_pixels},
\code{pass2_pixels}). When FT08 was run, \code{SpectralStemInferencer} overrode only the first. The second still
used cells of one token. \Cref{we:decoderwork:leak} counts what that did, and \cref{sec:decoderwork:leak} tells how
it was found and fixed.

\begin{worked}[label={we:decoderwork:leak}]{Half of each block taken from a pass that saw it}
A $4\times4$-token tile, block 2. Write $h(i,j)=1$ if pass 1 hides token $(i,j)$ (cells of 2 tokens, pass 1 hiding
the cells whose block row plus block column is odd) and $s(i,j)=1$ if the pixel's prediction is taken from pass 1
(cells of 1 token, $i+j$ odd). A token is predicted correctly when $s=h$: taken from pass 1 exactly when pass 1 hid
it.
\begin{center}\small\ttfamily
\begin{tabular}{cccc@{\qquad}cccc@{\qquad}cccc}
\multicolumn{4}{c}{$h$} & \multicolumn{4}{c}{$s$} & \multicolumn{4}{c}{leak ($s\ne h$)}\\
0&0&1&1 & 0&1&0&1 & .&x&x&.\\
0&0&1&1 & 1&0&1&0 & x&.&.&x\\
1&1&0&0 & 0&1&0&1 & x&.&.&x\\
1&1&0&0 & 1&0&1&0 & .&x&x&.\\
\end{tabular}
\end{center}
Count the leaks: two per row, $2+2+2+2=8$ of 16 tokens, $8/16=0.5$. The reason is local: inside any $2\times2$
block $h$ is constant, while $s$ takes each value on exactly 2 of the 4 tokens. So in every block, whichever way
the two patterns line up, 2 of 4 tokens agree and 2 disagree. \textbf{Half of each block's pixels were predicted by
the pass in which their block was visible.}
\end{worked}

\begin{sourcenote}[: checkerboard cell size is in tokens]
\code{InferenceConfig.checkerboard_cell_size} (\file{app/models/training/inference_config.py}) was described as
the ``Size of each checkerboard cell in pixels''. The code passes it to \code{checkerboard_token_mask} as
\code{cell_size}, documented there as ``checkerboard cell size in tokens'', and the inferencer compares it
with \code{mask_block}, also in tokens. It is a token count. The description was corrected to ``stage-1 tokens
(4 px each)'' on 2026-10-21 after the course found it.
\end{sourcenote}

\paragraph{What validation measures.} Both runs' per-epoch validation is ungated and uses single-token masks:
\code{compute_validation_loss} in \file{spectral_stem_mae_trainer.py} builds two complementary random token masks
itself, without \code{_build_prediction_mask} and without \code{gate_pixels}. So validation numbers compare
with FT04's, and for FT08 they are measured under a masking different from its training and its inference.

\section{FT07 and FT08: results}\label{sec:decoderwork:ft0708}

The pass criteria, fixed before launch, against FT04 + RM with cloud excluded on the real PRISMA bench:
(1) pAUC $\le$1\% at least $+0.005$ \emph{and} its paired-bootstrap 95\% interval above zero
(\cref{ch:bootstrap}); (2) objects caught at the top 1\% and 0.1\% not lower; (3) bench-v2.1 caught top 1\% not
lower by more than 0.5 points; (4) validation SAM within 0.1° of FT04 (3.88 / 4.11 / 5.46 for EnMAP / PRISMA /
AVIRIS-NG) and EnMAP IoU water / cloud within 0.01 (0.839 / 0.810).

\begin{table}[htbp]
\centering\small
\begin{tabular}{lccccc}
\toprule
 & pAUC & top 1\% / 0.1\% & TPR @ 0.1\% & ROC-AUC & v2.1 top 1\%\\
\midrule
FT04 & \textbf{0.170} & 46.0\% / 15.8\% & 0.065 & 0.852 & 36.2\%\\
FT07 & 0.165 & 46.1\% / 15.4\% & 0.062 & -- & 36.4\%\\
FT07 epoch 1 & 0.169 & 46.3\% / 15.5\% & 0.063 & -- & --\\
FT08 & 0.164 & 46.4\% / 16.1\% & 0.062 & 0.844 & 36.7\%\\
\midrule
\multicolumn{6}{l}{Paired bootstrap against FT04 (1000 draws), pAUC:}\\
\multicolumn{6}{l}{\quad FT07 $-0.0047\ [-0.0082,-0.0020]$;\quad FT08 $-0.0058\ [-0.0106,-0.0010]$}\\
\bottomrule
\end{tabular}
\caption{FT07 and FT08 detection, real PRISMA bench, cloud excluded on FT04's masks; ``top'' is objects caught.
Both best checkpoints are epoch 4. Source: \tagM{FT07-gated/README.md}.}
\label{tab:decoderwork:ft0708}
\end{table}

\begin{drift}
Run 2026-10-21. FT07 validation at epoch 4: SAM 3.90 / 4.16 / 5.47°, IoU 0.840 / 0.813. FT08 at epoch 4: SAM
3.92 / 4.17 / 5.33°, IoU 0.841 / 0.807; FT08's AVIRIS-NG SAM swung 5.78 (epoch 3), 5.33 (4), 5.60 (5), as large as
the 0.1° guard \tagM{FT07-gated/README.md}.
\end{drift}

Both failed criterion 1: the intervals lie wholly below zero, a real loss, small. Objects caught were a tie.
Reconstruction and masks held (criterion 4), so the cost is specific to detection. FT07's cost grew with
training: epoch 1 was $-0.001$, epoch 4 $-0.005$.

\paragraph{Why gating hit the wrong pixels.} Real targets are under 0.1\% of pixels, and many are not the
worst-reconstructed in a batch. The top 1\% of training error is mostly hard \emph{normal} pixels: field edges,
texture, glint, shoreline. Gating them leaves the model slightly worse exactly where RM's false alarms come
from, and does little to the targets.

\paragraph{Why blocks moved objects one way and pixels the other.} Objects caught went up a little for FT08
(unresolved) while pAUC fell. Hiding 8\,px blocks makes the model reconstruct each hidden pixel from further away
(\cref{sec:decoderwork:neighbours}), so ordinary texture is rebuilt less exactly: more background residual
everywhere, while targets gain little extra contrast.

\begin{keypoint}
Making the training ignore what it gets wrong does not make it ignore targets, because most of what it gets
wrong is ordinary hard background.
\end{keypoint}

\subsection{The leak in FT08's evaluation, and the re-run}\label{sec:decoderwork:leak}

The numbers above for FT08 were measured with the pass selection of \cref{we:decoderwork:leak}. On 2026-10-31,
while this book was being written, its authors found the mismatch, and it was fixed the same day (commit d025059).
\code{SpectralStemInferencer.__init__} now copies the model's block into the inference configuration before the
parent is built, so both checkerboard choices read the same number:
\begin{lstlisting}
def __init__(self, config):
    block = getattr(config.model_config_, "mask_block", 1)
    if block > config.checkerboard_cell_size:
        config = config.model_copy(update={"checkerboard_cell_size": block})
    super().__init__(config)
\end{lstlisting}
The old \code{_checkerboard_keep_mask} override is gone. A test pins it:
\code{test_block_inference_never_predicts_a_pixel_from_a_pass_that_saw_it}
(\file{tests/test_foundation_models/test_gated_training.py}) builds a small model with \code{mask_block=2}, adds
0.4 to every band of one $8\times8$-pixel block and checks that the predictions inside that block do not change.
If any of them came from a pass that saw the block, they would. The FT07 README records that the test fails
without the fix. FT04 and Heavy have \code{mask_block} 1 (the configuration default), so the \code{if} is false and
nothing changes for them.

FT08 was re-evaluated (\file{configs/eval/ft08_fixed.json}):
\begin{center}\small
\begin{tabular}{lccc}
\toprule
 & pAUC & caught top 1\% / 0.1\% & paired difference [95\% interval]\\
\midrule
FT08, leaky (2026-10-21) & 0.164 & 46.4\% / 16.1\% & $-0.0058\ [-0.0106,-0.0010]$ vs FT04\\
FT08, fixed (2026-10-31) & 0.1638 & 46.1\% / 15.4\% & $-0.0057\ [-0.0101,-0.0013]$ vs FT04\\
 & & & $-0.0010\ [-0.0047,+0.0028]$ vs FT07\\
\bottomrule
\end{tabular}

\smallskip
{\footnotesize Source: \tagM{research/experiments/indradhanu-finetune/FT07-gated/README.md} (last section).}
\end{center}
Against FT04 the clean FT08 is still a resolved loss; against FT07, the same gated loss without blocks, it is a
tie. Blocks added nothing on top of gating. \textbf{The verdict ``block masking does not help'' stands, now on a
measurement without the leak.}

Why did predicting half the pixels from a pass that could see them move pAUC by only $0.0002$? A likely reason,
not tested for FT08: \cref{sec:decoderwork:diagnosis} measures that FT04 reproduces pixels it can see barely
better than ones it must predict (set2 background 0.110 visible against 0.137 hidden). A model that cannot copy
what it sees gains little from seeing it.

\begin{lesson}[: one setting, read in two places]
A configuration value was read in two places and overridden in one. Nothing crashed, the numbers were
plausible, and the conclusion happened to survive; but nobody could know that without re-running the
evaluation. Two habits follow. When a model needs a non-default inference setting, change it where \emph{every}
reader of the setting sees it (here, the configuration), not in one of the methods that reads it. And write the
test that states the invariant directly (``no pixel is predicted by a pass that saw it''), so a future override
cannot quietly break it again. The leak was found by people reading the code to explain it, which is one argument
for writing explanations at all.
\end{lesson}

%======================================================================
\section{Is the decoder the bottleneck? A diagnosis without training}\label{sec:decoderwork:diagnosis}

Every change since FT04 had altered the loss, the masking or the score, never the shape of the network. The
user's prompt was: ``Maybe a bigger model is the answer'', ``Or maybe more layers of decoding'', ``I feel the
decoder maybe too un-expressive given the sophistication of the encoder.'' Before building anything, two
measurements on the laptop, with no training, asked \emph{where} FT04's error comes from
\tagM{decoder-capacity/README.md}.

\subsection{Where the parameters are}

\begin{center}
\begin{tabular}{lrl}
\toprule
component & parameters & role\\
\midrule
encoder & 3,331,072 & four SegFormer stages (\cref{ch:encoder})\\
decoder (reconstruction) & 2,164,736 & fuse at 1/4, then PixelShuffle(4) (\cref{ch:decoder})\\
segmentation head & 371,536 & water / cloud masks (\cref{ch:seghead})\\
\texttt{spectral\_in} & 14,721 & bands to 32 numbers per pixel (\cref{ch:specdecoder})\\
\texttt{spectral\_out} & 4,257 & 32 numbers per pixel to any band list\\
\bottomrule
\end{tabular}
\end{center}
Building the FT04-sized model on CPU confirms the encoder, decoder and stem figures \tagC; the script
and its output are in \file{research/textbook/tools/}. The decoder is 39\% of encoder plus decoder:
$2{,}164{,}736/5{,}495{,}808=0.39$. It is not small, so ``make the decoder bigger'' is not the obvious fix. The
question is not how many parameters it has but \emph{what it is able to draw}.

\begin{worked}[label={we:decoderwork:ft04dec}]{FT04's decoder, 2,164,736 parameters}
A \code{SegFormerDecoder} with widths $[32,64,160,256]$, \code{decoder_dim} 256 and 32 output channels (the
spectral latent), so the last convolution makes $32\times16=512$ sub-pixel channels.
\begin{align*}
\text{projections} &: (32+64+160+256)\times256 + 4\times256 = 512\times256+1{,}024\\
&\phantom{:}\ = 131{,}072+1{,}024 = 132{,}096\\
\text{fuse} &: 1{,}024\times256+256 = 262{,}144+256 = 262{,}400\\
\text{refine }3\times3 &: 256\times256\times9+256 = 589{,}824+256 = 590{,}080\\
\text{sub-pixel }3\times3 &: 256\times512\times9+512 = 1{,}179{,}648+512 = 1{,}180{,}160\\
\text{total} &: 132{,}096+262{,}400 = 394{,}496;\quad 394{,}496+590{,}080 = 984{,}576\\
&\phantom{:}\ 984{,}576+1{,}180{,}160 = 2{,}164{,}736
\end{align*}
Over half the decoder (1,180,160) is the one layer that emits a whole $4\times4$ block at once.
\end{worked}

\subsection{Test 1: the spectral floor}\label{sec:decoderwork:floor}

\code{WavelengthLinearDecoder} (\cref{ch:specdecoder}) makes band $n$ of a pixel as
\begin{equation}\label{eq:decoderwork:wld}
\hat x_n = \vect{U}(\lambda_n)\cdot\vect{z} + c(\lambda_n),
\end{equation}
with $\vect{z}$ the 32 numbers the decoder predicts for the pixel and $\vect{U}(\lambda_n)$ a 32-vector. Stack
the rows $\vect{U}(\lambda_n)$ into a matrix $\mat{U}$ of shape $N\times32$: every output spectrum is the offset
curve $\vect{c}$ plus $\mat{U}\vect{z}$, a mixture of the 32 fixed columns of $\mat{U}$.

\begin{definition}[Span and spectral floor]
The \emph{span}\index{span} of a set of vectors is the set of all their weighted sums. The \emph{spectral
floor}\index{spectral floor} of a pixel is the part of its (offset-removed) spectrum that lies outside the span of
the columns of $\mat{U}$: error that no choice of $\vect{z}$, and so no change upstream of \code{spectral_out},
can remove.
\end{definition}

The closest point of the span to a spectrum $\vect{x}$ is its \emph{orthogonal projection}\index{projection}
(\cref{ch:linalg}). For one curve $\vect{u}$ the projection is $a^\star\vect{u}$ with
$a^\star=(\vect{x}\cdot\vect{u})/(\vect{u}\cdot\vect{u})$.

\begin{worked}[label={we:decoderwork:floor}]{The floor with 3 bands and 1 curve}
Do this on paper. The model has 3 bands, one curve $\vect{u}=(1,1,1)$ and no offset: every reconstruction is a
flat spectrum $a(1,1,1)$. The real pixel is $\vect{x}=(2,1,3)$.
\begin{align*}
\vect{x}\cdot\vect{u} &= 2\times1+1\times1+3\times1 = 2+1+3 = 6\\
\vect{u}\cdot\vect{u} &= 1+1+1 = 3\\
a^\star &= 6/3 = 2;\quad\text{projection } = (2,2,2)\\
\text{floor} &= \vect{x}-\text{projection} = (2-2,\ 1-2,\ 3-2) = (0,-1,1)\\
\text{floor energy} &= 0^2+(-1)^2+1^2 = 0+1+1 = 2
\end{align*}
Now the model predicts $a=1.5$:
\begin{align*}
\text{residual} &= (1.5-2,\ 1.5-1,\ 1.5-3) = (-0.5,\ 0.5,\ -1.5)\\
\text{energy} &= 0.25+0.25+2.25 = 2.75\\
\text{floor share} &= 2/2.75 = 0.73
\end{align*}
73\% of this residual's energy is floor; the other 27\% is the model choosing 1.5 instead of 2.
\end{worked}

\paragraph{Why the share can never exceed 1.} Split the residual $\vect{e}=\xhat-\vect{x}$ into its part in the
span and its part at right angles to the span. Every reconstruction lies in the span, so the out-of-span part of
$\vect{e}$ is exactly minus the out-of-span part of $\vect{x}$, the floor $\vect{f}$. Write
$\vect{e}=\vect{s}-\vect{f}$ with $\vect{s}$ in the span and $\vect{f}$ orthogonal to it. Then
\begin{align*}
\norm{\vect{e}}^2 &= (\vect{s}-\vect{f})\cdot(\vect{s}-\vect{f})
  \intertext{expand the product:}
  &= \vect{s}\cdot\vect{s} - 2\,\vect{s}\cdot\vect{f} + \vect{f}\cdot\vect{f}
  \intertext{$\vect{s}\cdot\vect{f}=0$ because they are at right angles:}
  &= \norm{\vect{s}}^2+\norm{\vect{f}}^2 \;\ge\; \norm{\vect{f}}^2 .
\end{align*}
This is Pythagoras' theorem. Check it on \cref{we:decoderwork:floor}:
\[
\vect{s}=(1.5,1.5,1.5)-(2,2,2)=(-0.5,-0.5,-0.5),\qquad \norm{\vect s}^2=0.75,\qquad 0.75+2=2.75. \checkmark
\]

The script \file{spectral_floor.py} does this with 32 curves, using a QR decomposition: \code{Q, _ =
np.linalg.qr(U)} gives an orthonormal basis $\mat{Q}$ of the same span, so the projection is
$\mat{Q}(\mat{Q}\T\vect{x})$. It runs over every target pixel and 200,000 random background pixels per scene.

\begin{table}[htbp]
\centering\small
\begin{tabular}{lcccc}
\toprule
scene & background: actual / floor & floor share & targets: actual / floor & floor share\\
\midrule
set1 & 0.048 / 0.013 & 0.07 & 0.189 / 0.022 & 0.02\\
set2 & 0.122 / 0.017 & 0.02 & 0.204 / 0.025 & 0.02\\
set3 & 0.121 / 0.033 & 0.09 & 0.149 / 0.043 & 0.10\\
set4 & 0.115 / 0.015 & 0.03 & 0.177 / 0.025 & 0.03\\
\bottomrule
\end{tabular}
\caption{The spectral floor of FT04 on the PRISMA bench, normalised units. ``Actual / floor'' are mean absolute
errors; the share is of squared error (energy). Source: \tagM{decoder-capacity/README.md}.}
\label{tab:decoderwork:floor}
\end{table}

The floor is 2--10\% of the residual energy. Over 90\% of FT04's error is in which $\vect{z}$ the decoder
predicts. Test 1 clears \code{spectral_out}. (Mind the units: set1 background's floor is $0.013/0.048=0.27$ of
the error in mean-absolute terms but $0.07$ in energy terms; squaring shrinks small parts relative to large
ones.)

\subsection{Test 2: visible against hidden pixels}\label{sec:decoderwork:visible}

FT04's background error could come from two causes.
\begin{itemize}
\item \textbf{Unpredictability.} A hidden pixel's value is not fully fixed by its neighbours. No decoder can fix
  this, and it is also what makes targets detectable.
\item \textbf{Capacity.} The network cannot \emph{render} the detail even when it has it. A better architecture
  can fix this.
\end{itemize}
Run the model with \emph{no mask}, every token visible: unpredictability is gone, so what error is left is
capacity. Compare it with the error on hidden pixels from the product's two complementary passes. The script
\file{visible_vs_hidden.py} cuts 24 random 128\,px tiles per PRISMA scene, keeping only tiles at least 95\% valid.

\begin{table}[htbp]
\centering\small
\begin{tabular}{lcccc}
\toprule
scene & background visible & background hidden & targets visible & targets hidden\\
\midrule
set1 & 0.056 & 0.057 & 0.194 & 0.218\\
set2 & 0.110 & 0.137 & 0.161 & 0.220\\
set3 & 0.159 & 0.129 & 0.176 & 0.179\\
set4 & 0.135 & 0.139 & 0.169 & 0.191\\
\bottomrule
\end{tabular}
\caption{FT04's mean absolute error per band (normalised), with no mask (visible) and in the two-pass
inference (hidden). Source: \tagM{decoder-capacity/README.md}.}
\label{tab:decoderwork:visible}
\end{table}

On set1 FT04 copies the background it can see (0.056) no better than the background it must predict (0.057). On
set2 and set4 visible is only modestly better; on set3 it is \emph{worse}, which neither the README nor this book
can explain. A model limited by unpredictability would autoencode visible pixels almost exactly. \textbf{This one
is limited by what its architecture can render.}

\subsection{Why: the \texorpdfstring{$4\times4$}{4x4}-token bottleneck, counted}\label{sec:decoderwork:bottleneck}

Follow one stage-1 token. The stage-1 patch embedding is a convolution with kernel 4 and stride 4
(\cref{ch:patchembed}), so a token reads a $4\times4$ block of pixels with 32 channels each:
\begin{align*}
\text{numbers into the token} &= 4\times4\times32 = 512\\
\text{numbers out (stage-1 width)} &= 32\\
\text{compression} &= 512/32 = 16\text{ to }1;\quad\text{per pixel } 32/16 = 2 \text{ numbers}.
\end{align*}

\paragraph{A counting argument for a whole tile.} The decoder must produce $32\times128\times128=524{,}288$
numbers for a 128\,px tile (32 latent numbers per pixel). Everything it knows about the tile arrives through the
four encoder stages:
\begin{align*}
\text{stage 1} &: 32\times32\text{ tokens}\times32 = 32{,}768\\
\text{stage 2} &: 16\times16\times64 = 16{,}384\\
\text{stage 3} &: 8\times8\times160 = 10{,}240\\
\text{stage 4} &: 4\times4\times256 = 4{,}096\\
\text{total} &: 32{,}768+16{,}384+10{,}240+4{,}096 = 63{,}488\\
\text{per pixel} &: 63{,}488/16{,}384 = 3.875 \text{ numbers, against 32 to produce}.
\end{align*}
Only stage 1 is at the resolution of a $4\times4$ block, and it carries 2 numbers per pixel; the other stages are
coarser still. This is a count, not a proof: the network is nonlinear, and natural images are compressible, so
3.875 numbers per pixel can carry a lot. But the detail that does not fit through has no other route.

\paragraph{The linear case: rank.} If the path were linear, the argument would be exact. A linear map from 512
numbers to 32 and back to 512 is a product $\mat{D}\mat{E}$ with $\mat{E}$ of shape $32\times512$ and $\mat{D}$
of shape $512\times32$; its \emph{rank}\index{rank (of a matrix)} (the dimension of the set of outputs it can
produce) is at most 32. At least $512-32=480$ directions of variation inside a block are lost entirely.

\begin{worked}[label={we:decoderwork:rank}]{Two pixels through one number}
Squeeze two one-channel pixels $(a,b)$ through one number and back: encode $z=(a+b)/2$, decode $(z,z)$.
\begin{align*}
(a,b)=(3,1)&:\ z=(3+1)/2=2,\ \text{output }(2,2),\ \text{error }(2-3,\ 2-1)=(-1,\ 1)\\
(a,b)=(1,3)&:\ z=(1+3)/2=2,\ \text{output }(2,2),\ \text{error }(1,\ -1)\\
(a,b)=(2,2)&:\ z=2,\ \text{output }(2,2),\ \text{error }(0,0)
\end{align*}
Three different inputs, one output. The difference $a-b$ never reaches the decoder, so no decoder, however large,
can restore it, even though the input was \emph{visible}. That is test 2's finding in miniature: visible
pixels are copied no better than hidden ones.
\end{worked}

FT04's decoder then fuses the four stages at 1/4 resolution and emits each $4\times4$ block \emph{at once} from
one 256-wide vector: a $3\times3$ convolution to 512 channels and PixelShuffle(4) (\cref{ch:decoder}). Every
pixel gets its own output channel and its own weights, but all 16 are computed from the same token-level input.
\textbf{Nothing at full resolution ever comes back in.}

\begin{keypoint}
The decoder's limit is not its parameter count but its paths: sub-token detail that did not fit through 32
numbers per token is gone, for visible pixels as much as hidden ones.
\end{keypoint}

%======================================================================
\section{FT09: the progressive decoder}\label{sec:decoderwork:ft09}

The diagnosis asked for two things: upsampling in \emph{stages}, so detail is rendered gradually, and a
\emph{full-resolution skip}\index{skip connection!full-resolution}: a second input that bypasses the encoder
and arrives at the decoder at full size, in the style of U-Net \citep{ronneberger2015unet}. This is not the
residual connection inside a transformer block (\cref{ch:block}); it runs from the start of the network to the
end.

\subsection{Shapes, stage by stage}

The skip is $\vect{z}$, the stem output after BatchNorm: 32 numbers per pixel. \code{ProgressiveDecoder} in
\file{app/foundation_models/components/progressive_decoder.py}, with FT09's widths (\code{decoder_dim} 256,
\code{skip_channels} 32, \code{stage_dims} (192, 96)), for a 128\,px tile and batch $B$:

\begin{table}[htbp]
\centering\small
\begin{tabular}{lll}
\toprule
scale & operation & shape\\
\midrule
in & encoder stages F1, F2 & \shape{(B,32,32,32)}, \shape{(B,64,16,16)}\\
in & encoder stages F3, F4 & \shape{(B,160,8,8)}, \shape{(B,256,4,4)}\\
in & skip $\vect{z}$, hidden pixels zeroed & \shape{(B,32,128,128)}\\
1/4 & $1\times1$ projections to 256, bilinear up, concat & \shape{(B,1024,32,32)}\\
1/4 & $1\times1$ fuse + GELU (warm-started) & \shape{(B,256,32,32)}\\
1/2 & bilinear $\times2$ & \shape{(B,256,64,64)}\\
1/2 & concat \texttt{avg\_pool2d(skip, 2)} & \shape{(B,288,64,64)}\\
1/2 & \texttt{up\_half}: $3\times3$ conv, GELU, $3\times3$ conv, GELU & \shape{(B,192,64,64)}\\
1 & bilinear $\times2$ & \shape{(B,192,128,128)}\\
1 & concat skip & \shape{(B,224,128,128)}\\
1 & \texttt{up\_full}: $3\times3$ conv, GELU, $3\times3$ conv, GELU & \shape{(B,96,128,128)}\\
1 & \texttt{head}: $1\times1$ conv & \shape{(B,32,128,128)}\\
out & \texttt{spectral\_out}, unchanged & \shape{(B,N,128,128)}\\
\bottomrule
\end{tabular}
\caption{\texttt{ProgressiveDecoder.forward} for a 128\,px tile with FT09's widths.}
\label{tab:decoderwork:shapes}
\end{table}

\begin{figure}[htbp]
\centering
\begin{tikzpicture}[font=\small,box/.style={draw,rounded corners,minimum height=7mm,align=center},>=Stealth]
\node[box,fill=spat!15] (enc) at (0,0) {encoder\\F1..F4};
\node[box,fill=attn!12] (fuse) at (2.8,0) {project+fuse\\$256@32^2$};
\node[box,fill=attn!12] (half) at (5.8,0) {\texttt{up\_half}\\$192@64^2$};
\node[box,fill=attn!12] (full) at (8.8,0) {\texttt{up\_full}\\$96@128^2$};
\node[box,fill=spec!15] (head) at (11.4,0) {head\\$32@128^2$};
\node[box,fill=spec!15] (stem) at (0,-2.2) {stem $\vect z$\\$32@128^2$};
\node[box,fill=maskgrey!30] (zero) at (2.8,-2.2) {zero hidden\\pixels};
\node[box,fill=maskgrey!15] (pool) at (5.8,-2.2) {avg pool $2$\\$32@64^2$};
\draw[->] (enc)--(fuse); \draw[->] (fuse)--node[above]{$\times2$}(half); \draw[->] (half)--node[above]{$\times2$}(full); \draw[->] (full)--(head);
\draw[->] (stem)--(zero); \draw[->] (zero)--(pool); \draw[->] (pool)--(half);
\draw[->] (zero.east) -- ++(0.4,0) |- ($(zero.south)+(0,-0.4)$) -| (full.south);
\draw[->] (stem.north)--(enc.south);
\end{tikzpicture}
\caption{The progressive decoder. The stem output reaches the decoder twice: through the encoder (top) and
through the masked skip (bottom), at half and full resolution. Channel counts are written as ``channels @ size''.}
\label{fig:decoderwork:progressive}
\end{figure}

\subsection{The parameter count by hand}

\begin{worked}[label={we:decoderwork:prog}]{The progressive decoder's 1,504,096 parameters}
Every convolution has $c_\text{in}\times c_\text{out}\times k\times k+c_\text{out}$ parameters. Do it on paper.
\begin{align*}
\text{projections} &: (32+64+160+256)\times256 + 4\times256 = 131{,}072+1{,}024 = 132{,}096\\
\text{fuse} &: 1{,}024\times256+256 = 262{,}144+256 = 262{,}400\\
\texttt{up\_half}\ 1 &: 288\times192\times9+192 = 497{,}664+192 = 497{,}856\\
\texttt{up\_half}\ 2 &: 192\times192\times9+192 = 331{,}776+192 = 331{,}968\\
\texttt{up\_full}\ 1 &: 224\times96\times9+96 = 193{,}536+96 = 193{,}632\\
\texttt{up\_full}\ 2 &: 96\times96\times9+96 = 82{,}944+96 = 83{,}040\\
\texttt{head} &: 96\times32+32 = 3{,}072+32 = 3{,}104\\[3pt]
\text{running total} &: 132{,}096+262{,}400 = 394{,}496\\
&\phantom{:}\ 394{,}496+497{,}856 = 892{,}352\\
&\phantom{:}\ 892{,}352+331{,}968 = 1{,}224{,}320\\
&\phantom{:}\ 1{,}224{,}320+193{,}632 = 1{,}417{,}952\\
&\phantom{:}\ 1{,}417{,}952+83{,}040 = 1{,}500{,}992\\
&\phantom{:}\ 1{,}500{,}992+3{,}104 = 1{,}504{,}096
\end{align*}
\end{worked}

\begin{sourcenote}[: 1,504,096 verified]
Building \code{ProgressiveDecoder([32, 64, 160, 256])} on CPU gives 1,504,096 \tagC
(\file{research/textbook/tools/compute_39-41.py}); the source comment reads ``1.50 M with FT04 widths (FT04's
decoder: 2.16 M)''. \textbf{The new decoder has fewer parameters than the old one}: it is wired differently, not
bigger.
\end{sourcenote}

It does far more arithmetic, though. \texttt{up\_full}'s first convolution runs its 193,536 multiply-adds at each
of $128\times128=16{,}384$ pixels: $193{,}536\times16{,}384\approx3.17\times10^9$ per tile. FT04's two refine
convolutions together at $32\times32$: $(589{,}824+1{,}179{,}648)\times1{,}024=1{,}769{,}472\times1{,}024\approx
1.81\times10^9$. These are hand counts, not timings \tagC.

\subsection{The skip must not carry hidden pixels}

This is the trap in the design. The masked autoencoder works because a hidden pixel's value never reaches its
own reconstruction (\cref{ch:recon}). A skip that bypasses the encoder would undo that: hand the decoder the real
pixel and it will copy it, the residual vanishes, and every target with it. So \code{_decode} in
\file{spectral_stem_seg_former_mae.py} zeroes them:
\begin{lstlisting}
visible = keep_mask.reshape(B, 1, H // 4, W // 4).to(z.dtype)    # (B, 1, 32, 32): 1 = visible token
z = z * F.interpolate(visible, scale_factor=4, mode="nearest")   # (B, 32, 128, 128): token's 16 pixels
return self.decoder(features, z)
\end{lstlisting}
Nearest-neighbour interpolation by 4 copies each token's 0 or 1 to its 16 pixels. The half-resolution skip is
\code{avg_pool2d(skip, 2)}: since 2 divides 4 and the $2\times2$ windows line up with token borders, each average
stays inside one token, so a hidden token is 0 at both resolutions. Neighbouring \emph{visible} pixels still reach
a hidden pixel through the $3\times3$ convolutions; the docstring calls that context, as in the encoder.

The test \code{test_hidden_pixels_never_reach_the_output} builds a small model and a 64\,px tile ($16\times16=256$
tokens), hides tokens 0, 17, 100 and 255, makes a second input in which \emph{only those tokens' pixels} are
replaced with fresh random values, and asserts the two reconstructions agree everywhere to $10^{-6}$. Token $t$
is at row $\lfloor t/16\rfloor$, column $t \bmod 16$: token 17 is $(1,1)$, pixels rows 4--7, columns 4--7; token
100 is $(6,4)$, rows 24--27, columns 16--19. Its partner test shows the skip is live: adding 0.5 to one visible
pixel moves that pixel's own reconstruction by more than $10^{-4}$.

\subsection{Warm start and two-phase training}

FT09 starts from FT04's \code{best.pt}. The new decoder kept the old layer names \code{linear_projections} and
\code{fuse_conv}, so those load; the old \code{refine.*} is dropped. How much of the decoder is warm?
$132{,}096+262{,}400=394{,}496$ of 1,504,096, that is $394{,}496/1{,}504{,}096=0.26$. \textbf{Three quarters of the
new decoder starts random.}

The first attempt trained everything jointly from the start (lr $10^{-4}$, planned 8 epochs) and was stopped
after epoch 3, still 0.5--1° SAM behind FT04: a mostly random decoder had $3\times30{,}000=90{,}000$ samples to
learn from, against the millions behind FT04's, on a learning rate already decaying, and its early gradients
pushed a good encoder around to suit a bad decoder. The fix split training in two \tagM{FT09-decoder/README.md}:

\begin{center}\small
\begin{tabular}{llllll}
\toprule
run & mode & what trains & lr & epochs & starts from\\
\midrule
FT09a & \texttt{"decoder"} & \texttt{decoder.*} only & $10^{-3}$ & 3 & FT04\\
FT09b & \texttt{"joint"} & everything & $5\times10^{-5}$ & 5 & FT09a\\
FT09c & \texttt{"joint"} & everything & $3\times10^{-5}$ & 5 & FT09b\\
\bottomrule
\end{tabular}
\end{center}
With \code{spectral_out} frozen in FT09a, the decoder is in effect learning to emit FT04's own code $\vect{z}$.
FT09a reached FT04's reconstruction, within 0.1° on all three sensors (3.98 / 4.14 / 5.50° against 3.88 / 4.11 /
5.46°), in 3 epochs.

\subsection{Did it do what it was built for? Contrast}

\begin{table}[htbp]
\centering\footnotesize\setlength{\tabcolsep}{3pt}
\begin{tabular}{lcccc}
\toprule
 & set1 & set2 & set3 & set4\\
\midrule
hidden background FT04 $\to$ FT09a & 0.057 $\to$ \textbf{0.051} & 0.137 $\to$ \textbf{0.125} & 0.129 $\to$ \textbf{0.117} & 0.139 $\to$ \textbf{0.110}\\
hidden targets FT04 $\to$ FT09a & 0.218 $\to$ 0.204 & 0.220 $\to$ 0.203 & 0.179 $\to$ 0.171 & 0.191 $\to$ 0.162\\
target / background & 3.82 $\to$ \textbf{4.00} & 1.61 $\to$ \textbf{1.62} & 1.39 $\to$ \textbf{1.46} & 1.37 $\to$ \textbf{1.47}\\
visible background & 0.056 $\to$ 0.176 & 0.110 $\to$ 0.120 & 0.159 $\to$ 0.120 & 0.135 $\to$ 0.115\\
\bottomrule
\end{tabular}
\caption{The mechanism check on FT09a, mean absolute error per band, normalised. Source:
\tagM{FT09-decoder/README.md}.}
\label{tab:decoderwork:contrast}
\end{table}

\begin{worked}[label={we:decoderwork:contrast}]{Set1's contrast change}
\begin{align*}
\text{background fall} &: 0.057-0.051 = 0.006;\quad 0.006/0.057 = 0.105\ (10.5\%)\\
\text{target fall} &: 0.218-0.204 = 0.014;\quad 0.014/0.218 = 0.064\ (6.4\%)\\
\text{new ratio} &: 0.204/0.051 = 4.00
\end{align*}
Background fell by a larger fraction than targets, so the ratio rose.
\end{worked}

Hidden background error fell more than target error on all four scenes, so contrast rose on all four: the
intended effect, and the first model change in the series to move it.

\begin{sourcenote}[: ``6--21\%'' or ``9--21\%'']
Course part 67 records that the FT09 README said hidden background error ``falls 6-21\%'', while its own table
gives set1 10.5\%, set2 $0.012/0.137=8.8\%$, set3 $0.012/0.129=9.3\%$, set4 $0.029/0.139=20.9\%$: 9--21\%. The
README now reads ``falls 9-21\%'' (\file{research/experiments/indradhanu-finetune/FT09-decoder/README.md}); the
course note predates the correction.
\end{sourcenote}

\paragraph{Criterion 5, stated honestly.} Before launch the README fixed a mechanism criterion: visible-pixel
background error well below FT04's on all four scenes. The last row of \cref{tab:decoderwork:contrast} is that
test, and it \textbf{failed}: set1 got much worse (0.056 to 0.176). Afterwards the README recorded a caveat: the
progressive decoder only ever trained with about 65\% of its skip zeroed, so a fully unmasked input is outside
anything it has seen. That argument is reasonable, but the yardstick was judged wrong \emph{after} the result came
in, which is exactly what fixing criteria in advance is meant to prevent. Treat the mechanism as supported, not
proven.

\subsection{FT09b: detection, and the bootstrap}

\begin{table}[htbp]
\centering\small
\begin{tabular}{lccccc}
\toprule
 & pAUC & top 1\% / 0.1\% & TPR @ 0.1\% & ROC-AUC & v2.1 top 1\%\\
\midrule
FT04 & 0.170 & 46.0\% / \textbf{15.8\%} & \textbf{0.065} & 0.852 & \textbf{36.2\%}\\
FT09a & 0.169 & 46.1\% / 15.6\% & 0.061 & 0.856 & 35.8\%\\
FT09b & \textbf{0.176} & \textbf{47.1\%} / 15.4\% & 0.063 & \textbf{0.854} & 36.1\%\\
FT09c & 0.172 & 46.4\% / 15.0\% & -- & -- & 36.3\%\\
\bottomrule
\end{tabular}
\caption{FT09 detection, real PRISMA bench, cloud excluded on FT04's masks. Source: \tagM{FT09-decoder/README.md}.}
\label{tab:decoderwork:ft09}
\end{table}

The bench pAUC is the mean over the four scenes. Check it from the unrounded per-scene values in
\file{eval_per_scene.csv}:
\begin{align*}
\text{FT04} &: 0.24603+0.12817+0.08428+0.21986 = 0.67834;\quad 0.67834/4 = 0.16959\\
\text{FT09b} &: 0.26010+0.13061+0.08301+0.22873 = 0.70245;\quad 0.70245/4 = 0.17561\\
\text{change} &: 0.17561-0.16959 = 0.00602
\end{align*}
which is the bootstrap's point value $+0.0060$.

The paired object bootstrap of \cref{ch:bootstrap} resamples the bench's 909 objects with replacement 1000 times
and recomputes the difference each time; the 95\% interval runs from the 2.5th to the 97.5th percentile of those
differences. For FT09b $-$ FT04 \tagM{FT09b-decoder-joint/bootstrap.txt}:
\[
\text{pAUC: } +0.0060\ [+0.0002,\ +0.0118],\quad P(B>A)=0.976 .
\]
The interval clears zero, by $0.0002$. Caught at the top 1\%: $+1.1$ points $[-0.3,+2.4]$, a tie. Caught at the
top 0.1\%: $-0.4$ points $[-1.4,+0.6]$, a tie in the wrong direction: FT04 caught $57+33+13+41=144$ objects and
FT09b $56+34+9+41=140$, and $4/909=0.44\%$; the whole loss is set3, 13 to 9.

\begin{worked}[label={we:decoderwork:interval}]{Reading an interval as a margin}
The interval is centred near the point value: $0.0060-0.0002=0.0058$ below and $0.0118-0.0060=0.0058$ above.
Its half-width $0.0058$ is the bench's resolution for this pair. A gain of $+0.0060$ is only
$0.0060/0.0058=1.03$ half-widths from zero: the smallest margin the rule (``interval above zero'') allows.
Compare cloud exclusion, $+0.0074\ [+0.0047,+0.0101]$: half-width $0.0027$, margin $0.0074/0.0027=2.7$.
\end{worked}

So FT09b met criterion 1 by $+0.0002$, failed criterion 2 at the top 0.1\% (unresolved), met 3 ($-0.1$ points) and
4, and failed 5 as written. A candidate, not a release. FT09a alone, with higher contrast already, \emph{tied}
FT04 ($-0.0006\ [-0.0056,+0.0043]$): the joint phase, with the encoder adapting to the new decoder, did the work.

\subsection{FT09c: the gain does not hold}

FT09b's best checkpoint was its last epoch, so FT09c continued it for five epochs at $3\times10^{-5}$. Its
reconstruction got better (AVIRIS-NG SAM 5.41 to 5.34°, better than FT04) and its detection worse: pAUC 0.172, a
tie with FT04 ($+0.0020\ [-0.0026,+0.0064]$), and a \emph{resolved} loss against FT09b,
$-0.0040\ [-0.0076,-0.0010]$ \tagM{FT09-decoder/README.md}.

\begin{keypoint}
A gain must hold across neighbouring checkpoints: the object bootstrap resamples objects, not training runs, so
it cannot see checkpoint-to-checkpoint variation, which here was as large as the gain itself.
\end{keypoint}

The progressive decoder was not adopted. Its mechanism (higher hidden-pixel contrast) was real; its effect on
detection was within checkpoint noise. And again, as with FT06: better reconstruction is not better detection.

\begin{drift}
FT09 runs 2026-10-22, FT09c evaluated 2026-10-23 (it finished unevaluated and the GPU sat idle about 10 hours
because no watcher was set). Per-scene pAUC FT09c: 0.254 / 0.131 / 0.082 / 0.219
\tagM{FT09-decoder/README.md}.
\end{drift}

\section{The cumulative lesson since FT04}\label{sec:decoderwork:cumulative}

\begin{table}[htbp]
\centering\small
\begin{tabular}{llcl}
\toprule
change tried on FT04 & chapter & pAUC change [95\% interval] & verdict\\
\midrule
cloud exclusion & \ref{ch:bootstrap} & $+0.0074\ [+0.0047,+0.0101]$ & real gain, adopted\\
FT06 residual$/\sqrt b$ & \ref{ch:uncertainty} & $-0.0067\ [-0.0174,+0.0037]$ & tie\\
gated scene adaptation, $N=50$ & \ref{ch:scoringgains} & $-0.0091\ [-0.0176,-0.0028]$ & real loss\\
FT07 gated loss & this & $-0.0047\ [-0.0082,-0.0020]$ & real loss, small\\
FT08 gated loss + blocks & this & $-0.0057\ [-0.0101,-0.0013]$ & real loss, small\\
FT09a decoder only & this & $-0.0006\ [-0.0056,+0.0043]$ & tie\\
FT09b decoder, joint & this & $+0.0060\ [+0.0002,+0.0118]$ & gain by a hair\\
FT09c five more epochs & this & $+0.0020\ [-0.0026,+0.0064]$ & tie: FT09b did not hold\\
\bottomrule
\end{tabular}
\caption{Every change tried on FT04 with a paired bootstrap interval. Sources: \tagM{bench-resolution/README.md},
\tagM{FT07-gated/README.md}, \tagM{FT09-decoder/README.md}. FT06 residual$/b$ (0.170 to 0.125) was not
bootstrapped. FT08's row is the 2026-10-31 re-run without the evaluation leak (\cref{sec:decoderwork:leak}); the
leaky run read $-0.0058\ [-0.0106,-0.0010]$.}
\label{tab:decoderwork:cumulative}
\end{table}

\begin{figure}[htbp]
\centering
\begin{tikzpicture}
\begin{axis}[width=0.9\linewidth,height=7.5cm,xmin=-0.02,xmax=0.014,ymin=0.4,ymax=8.6,
  ytick={1,...,8},yticklabels={FT09c,FT09b,FT09a,FT08,FT07,gated SA,FT06 sqrt,cloud excl.},
  xlabel={pAUC change against FT04},xtick={-0.02,-0.015,-0.01,-0.005,0,0.005,0.01},
  scaled x ticks=false,xticklabel style={/pgf/number format/fixed,/pgf/number format/precision=3,font=\scriptsize},
  grid=major,grid style={gray!15}]
\draw[thick,maskgrey] (axis cs:0,0.4) -- (axis cs:0,8.6);
\addplot[only marks,mark=*,loss,error bars/.cd,x dir=both,x explicit]
 coordinates {
 (0.0074,8) += (0.0027,0) -= (0.0027,0)
 (-0.0067,7) += (0.0104,0) -= (0.0107,0)
 (-0.0091,6) += (0.0063,0) -= (0.0085,0)
 (-0.0047,5) += (0.0027,0) -= (0.0035,0)
 (-0.0057,4) += (0.0044,0) -= (0.0044,0)
 (-0.0006,3) += (0.0049,0) -= (0.0050,0)
 (0.0060,2) += (0.0058,0) -= (0.0058,0)
 (0.0020,1) += (0.0044,0) -= (0.0046,0)
 };
\end{axis}
\end{tikzpicture}
\caption{\Cref{tab:decoderwork:cumulative} drawn to scale: point values with their 95\% paired-bootstrap
intervals. Only cloud exclusion lies wholly right of zero; FT09b touches it, and FT09c falls back.}
\label{fig:decoderwork:intervals}
\end{figure}

Read down the last column. Since FT04, every change except cloud exclusion has been inside the bench's noise or a
small, consistent cost, and the one exception from the network's shape (FT09b) did not survive further training.
None of these changes altered how much the \emph{encoder} can hold: its widths stayed 32, 64, 160, 256. That is
the question \cref{ch:heavy} takes up.

\begin{confusions}
\confusion{Top 1\% gated means 1\% of the loss is gone.}{The halo multiplies it: one isolated pixel with
\texttt{gate\_halo = 2} becomes a $5\times5$ square of 25. And the share of loss \emph{value} removed exceeds the
share of pixels, because the gated pixels are the large ones (\cref{we:decoderwork:share}).}
\confusion{FT08 hides more of the tile than FT04.}{Almost the same share, 664 against 665 tokens on a 128\,px
tile, in bigger pieces. What changes is the visible context: 1.765 against 2.807 expected visible neighbours.}
\confusion{The FT08 leak inflated its score, so FT08 was really much worse.}{The clean re-run reads 0.1638
against the leaky 0.164, well inside what the bench resolves. The leak was real; its effect on this model was
tiny, and the verdict is the same.}
\confusion{FT08's validation score beat FT04's, so FT08 reconstructs better.}{Validation uses single-token masks,
not FT08's blocks, and SAM swings by about the guard from epoch to epoch.}
\confusion{The decoder has 2.16 M parameters, so it cannot be under-powered.}{Parameters are not the issue; paths
are. Its parameters all work at 1/4 resolution on token-level vectors (\cref{we:decoderwork:rank}).}
\confusion{A floor share of 0.07 means the floor is 7\% of the mean absolute error.}{The share is of squared
error (energy). In mean-absolute terms set1 background's floor is 0.27 of the error.}
\confusion{The progressive decoder is bigger, so the gain came from parameters.}{It is smaller, 1.50 M against
2.16 M; it does more arithmetic because it works at full resolution.}
\confusion{The skip leaks hidden pixels.}{Hidden tokens are zeroed at full and half resolution, and a test
checks it to $10^{-6}$. Visible neighbours do reach a hidden pixel through $3\times3$ convolutions; that is context.}
\confusion{FT09b cleared the bootstrap, so it is a real improvement.}{It cleared the \emph{object} bootstrap at
one checkpoint, by 0.0002. Five more epochs gave a resolved loss against it. The bootstrap cannot see checkpoint
noise.}
\end{confusions}

\begin{checkyourself}
\item\label{ex:decoderwork:1} A group has 200 loss pixels and \texttt{gate\_pct = 1}. How many are gated before
the halo? With \texttt{gate\_halo = 1}, how many loss pixels at most are removed, if the gated pixels are far apart
and away from edges? What share of L1 gradient terms is that?
\item\label{ex:decoderwork:2} A 256\,px tile with \texttt{mask\_block = 2} and ratio 0.65: how many blocks, how many
hidden, how many tokens hidden? Then compute the expected number of visible neighbours of a hidden token, for
blocks and for single tokens (where $4{,}096\times0.65$ is truncated).
\item\label{ex:decoderwork:3} With one curve $\vect{u}=(1,1,1)$ and pixel $\vect{x}=(4,1,1)$: find the projection,
the floor and its energy. If the model predicts $(3,3,3)$, what is the floor share? Verify Pythagoras.
\item\label{ex:decoderwork:4} Count the parameters of \texttt{up\_half}'s first convolution with
\texttt{stage\_dims = (160, 96)}, and the whole progressive decoder with that change (the second \texttt{up\_half}
convolution and \texttt{up\_full}'s first also change).
\item\label{ex:decoderwork:5} A colleague proposes stage-1 width 64 for FT04. What is the per-token compression,
and how many numbers per pixel would the four encoder stages then deliver for a 128\,px tile (other stages
unchanged)?
\item\label{ex:decoderwork:6} Set4's hidden background fell 0.139 to 0.110 and targets 0.191 to 0.162. Compute
each fall as a percentage and the new ratio. Then explain in two sentences why FT09b is not shipped.
\item\label{ex:decoderwork:7} Redo \cref{we:decoderwork:leak} on a $6\times6$-token tile. How many blocks, how
many tokens leak per block, and how many in all? Then explain why the answer does not depend on which pass hides
which blocks.
\end{checkyourself}

\nextchapter{ch:heavy}
